\documentclass{llncs}
\usepackage[utf8]{inputenc}


%% Save the class definition of \subparagraph
\let\llncssubparagraph\subparagraph
%% Provide a definition to \subparagraph to keep titlesec happy
\let\subparagraph\paragraph
%% Load titlesec
\usepackage[compact]{titlesec}
%% Revert \subparagraph to the llncs definition
\let\subparagraph\llncssubparagraph
% Add QED symbol at the end of proofs
\let\proof\relax\let\endproof\relax
\usepackage{amsthm}

%\titleclass{\ex}{straight}[\section]
%\newcounter{ex}

\usepackage{enumitem}

\usepackage[
n,
operators,
advantage,
sets,
adversary,
landau,
probability,
notions,
logic,
ff,
mm,
primitives,
events,
complexity,
asymptotics,
keys,oracles]{cryptocode}

\newenvironment{sol}{\par\noindent\rule{\textwidth}{0.4pt}
\paragraph{Solution:}}{\par\noindent\rule{\textwidth}{0.4pt}
}


\usepackage{wrapfig}
\usepackage{framed}
\usepackage[margin=2.5cm]{geometry}
\usepackage{hyperref}
\usepackage{zref-clever}
\newcommand{\cref}[1]{\zcref{#1}}

\newcommand{\half}{{\textstyle\frac{1}{2}}}

\newcommand{\SD}[2]{\Delta(#1 \: ;  #2)}

\newcommand{\cA}{{\mathcal A}}
\newcommand{\cO}{{\mathcal O}}
\newcommand{\cB}{{\mathcal B}}
\newcommand{\cC}{{\mathcal C}}
\newcommand{\cD}{{\mathcal D}}
\newcommand{\cP}{{\mathcal P}}
\newcommand{\cR}{{\mathcal R}}
\newcommand{\cQ}{{\mathcal Q}}
\newcommand{\cF}{{\mathcal F}}
\newcommand{\cG}{{\mathcal G}}
\newcommand{\cE}{{\mathcal E}}
\newcommand{\cH}{{\mathcal H}}
\newcommand{\cJ}{{\mathcal J}}
\newcommand{\cM}{{\mathcal M}}
\newcommand{\cV}{{\mathcal V}}
\newcommand{\cL}{{\mathcal L}}
\newcommand{\cS}{{\mathcal S}}
\newcommand{\cX}{{\mathcal X}}
\newcommand{\cY}{{\mathcal Y}}
\newcommand{\cZ}{{\mathcal Z}}
\newcommand{\cK}{{\mathcal K}}
\newcommand{\cI}{{\mathcal I}}
\newcommand{\cT}{{\mathcal T}}
\newcommand{\cW}{{\mathcal W}}
\newcommand{\sP}{{\mathsf{P}}}
\newcommand{\sQ}{{\mathsf{Q}}}
\newcommand{\sA}{{\mathsf{A}}}
\newcommand{\sB}{{\mathsf{B}}}
\newcommand{\sC}{{\mathsf{C}}}
\newcommand{\sD}{{\mathsf{D}}}


\usepackage{fullpage}


%\titleformat{\ex}{\normalsize\bfseries}{}{0em}{Exercise \theex:~}
%\titlespacing*{\ex}{0pt}{3.25ex plus 1ex minus .2ex}{1.5ex plus .2ex}

\renewcommand{\thesubsection}{Exercise \arabic{section}.\arabic{subsection}}
\zcRefTypeSetup{subsection}{
	Name-sg = Exercise ,
	name-sg = exercise ,
	Name-pl = Exercises ,
	name-pl = exercises ,
}
\renewcommand{\verify}{\mathsf{Verify}}
\usepackage{graphicx}
\usepackage{xcolor}

\newcommand{\DES}{{\mathsf{DES}}}
\newcommand{\negate}[1]{\overline{#1}}
\newcommand{\polyn}{\mathsf{poly}}
\newcommand{\Tag}{\mathsf{Tag}}
\newcommand{\advntg}{\mathsf{Adv}}
\newcommand{\wins}{\mathsf{wins}}

\newcommand{\PRNG}{\mathsf{PRNG}}

\newcommand{\suffix}{\mathsf{suffix}}
\newcommand{\CRT}{\mathsf{CRT}}
\newcommand{\puzzle}{\mathsf{puzzle}}
\newcommand{\setup}{\mathsf{Setup}}
\newcommand{\commit}{\mathsf{Commit}}
\newcommand{\open}{\mathsf{Open}}

\newcommand{\acc}{\mathsf{accept}}
\newcommand{\rej}{\mathsf{reject}}

%\newcommand{\floor}[1]{\lfloor #1 \rfloor}
\newcommand{\bins}{\bin^*}
\newcommand{\conc}{\: || \:}
\renewcommand{\pp}{\mathsf{pp}}
\usepackage{mdframed}

\mdfdefinestyle{protstyle}{%
        % userdefinedwidth=.5\linewidth,
        align=center,
	linecolor=black,linewidth=1.5pt,%
	frametitlerule=true,%
	frametitlebackgroundcolor=gray!20,
	innertopmargin=\topskip,
        skipabove=12pt,
}

\newenvironment{prot}[1]{\begin{mdframed}[style=protstyle,frametitle={#1}]}{\end{mdframed}\vspace{.3cm}}
\newcommand{\Exp}{\mathsf{Exp}}
\newcommand{\SSadv}{\mathsf{SSadv}}
\newcommand{\ct}{\mathsf{ct}}
\newcommand{\encap}{\mathsf{Encaps}}
\newcommand{\decap}{\mathsf{Decaps}}
\newcommand{\partdec}{\mathsf{partdec}}
\newcommand{\findec}{\mathsf{findec}}

\providecommand{\ifhideproofs}{\iffalse}
\ifhideproofs
\usepackage{environ}
\NewEnviron{hide}{}
\let\proof\hide
\let\endproof\endhide
\fi

\newcommand{\VRF}{\mathsf{VRF}}
\newcommand{\wuf}{\mathsf{WUF}}
\newcommand{\euf}{\mathsf{EUF}}
\newcommand{\qsdh}{\mathsf{qSDH}}
\newcommand{\G}{\mathbb{G}}
\newcommand{\com}{\mathsf{com}}
\newcommand{\F}{\mathbb{F}}
\newcommand{\Z}{\mathbb{Z}}
\newcommand{\R}{\mathbb{R}}


\title{Cryptography Exercises}
\subtitle{\textcolor{red}{under construction\\ do not disitribute\\   \scriptsize{\emph{last change \today}}}}
\author{Stefan Dziembowski}
\date{\today}


\institute{University of Warsaw}

\begin{document}

\maketitle




\section{Statistical distance and information-theoretic encryption}

Let $A$ and $B$ be two random variables taking values in a finite nonempty set $\cX$. Define the \emph{statistical distance\footnote{This is also known as the \emph{total variation distance}, especially in the quantum information literature} between $A$ and $B$} as
\[
\SD{A}{B} := \half \cdot \sum_{a \in \cX} \abs{\prob{A=a} - \prob{B=a}},
\]
and \emph{statistical distance of $A$ from uniformity} as
\[
d(A) := \SD{A}{U_\cX},
\]
where $U_\cX$ is uniform over $\cX$. The reference uniform distribution is always on the declared set; for a tuple, it is on the Cartesian product of the declared sets for its coordinates. Analogously, $\Delta$ can be defined as a distance between two \emph{probability distributions $\alpha,\beta:\cX \rightarrow [0,1]$} as
\[
\SD{\alpha}{\beta} := \SD{A}{B},
\]
where $A$ and $B$ are random variables (taking values from $\cX$) distributed according to $\alpha$ and $\beta$ (respectively). A similar convention applies to the distance $d$ from uniformity.

Randomized machines use private randomness independent of all external random choices.



\subsection{An alternative definition of statistical distance}
Show that for every $A$ and $B$ we have
\[
\SD{A}{B} = \sum_{a \in \cX:\, \prob{A=a} > \prob{B=a}} \left(\prob{A=a} - \prob{B=a}\right).
\]

\subsection{Statistical distance as a metric}
Let $\Pi$ be the set of all probability distributions over a finite nonempty set $\cX$. Prove that $\Delta$ is a \emph{metric} on this set, i.e., it satisfies the following axioms:
	\begin{itemize}
		\item \emph{non-negativity:} for every $\alpha,\beta \in \Pi$ we have $\SD{\alpha}{\beta} \geq 0$,
		\item \emph{identity of indiscernibles:} for every $\alpha,\beta \in \Pi$ we have that $\SD{\alpha}{\beta} = 0$ if and only if $\alpha = \beta$,
		\item \emph{symmetry:} for every $\alpha,\beta \in \Pi$ we have that $\SD{\alpha}{\beta} = \SD{\beta}{\alpha}$, and
		\item \emph{triangle inequality:} for every $\alpha,\beta,\gamma \in \Pi$ we have that $\SD{\alpha}{\gamma} \leq \SD{\alpha}{\beta} + \SD{\beta}{\gamma}$.
	\end{itemize}

\subsection{Statistical distance as distinguishing advantage}
Let $A$ and $B$ be two random variables taking values from some finite set $\cX$. For a (possibly randomized) machine $\cD$ that takes inputs from $\cX$ and outputs a bit, define its \emph{distinguishing advantage} between $A$ and $B$ as
\[
    \advntg^{\mathrm{dist}}_{\cD}(A,B) := \abs{\prob{\cD(A) = 1} - \prob{\cD(B) = 1}}.
\]
Prove that for every computationally unbounded machine $\cD$ we have
\begin{equation}\label{eq:AB}
    \advntg^{\mathrm{dist}}_{\cD}(A,B) \leq \SD{A}{B}.
\end{equation}
For every $A$ and $B$, construct $\cD$ that achieves equality in Eq.~(\ref{eq:AB}). Show that the maximum is unchanged if the absolute value is omitted, and hence
\[
    \SD{A}{B} = \max_{\cD}\left(\prob{\cD(A) = 1} - \prob{\cD(B) = 1}\right),
\]
where the maximum is over all computationally unbounded machines that output a bit.

\subsection{Statistical distance and guessing advantage}\label{ex:game}
Let $\alpha_0$ and $\alpha_1$ be two distributions over a finite nonempty set $\cX$. Consider the following game played by a computationally unbounded, possibly randomized machine $\cM$ (that knows $\alpha_0$ and $\alpha_1$):
\begin{enumerate}
		\item A uniformly random bit $b \leftarrow \bin$ is chosen.
		\item A value $A$ is sampled according to distribution $\alpha_b$ and sent to $\cM$.
		\item $\cM$ receives $A$ and produces output $\hat b \in \bin$.
\end{enumerate}
We say that $\cM$ \emph{won the game} if $\hat b = b$. Define its \emph{guessing advantage} as
\[
    \advntg^{\mathrm{guess}}_{\cM} := \abs{\prob{\cM\mbox{\ wins the game}} - \half}.
\]
Show that
\begin{equation}\label{eq:a}
	\forall_\cM \prob{\cM\mbox{\ wins the game}} \leq \frac{1 + \SD{\alpha_0}{\alpha_1}}{2}.
\end{equation}
For every $\alpha_0$ and $\alpha_1$, construct $\cM$ that achieves equality in Eq.~(\ref{eq:a}). By also considering the machine that complements its output, deduce that
\[
    \max_{\cM}\advntg^{\mathrm{guess}}_{\cM} = \half\SD{\alpha_0}{\alpha_1}.
\]

\subsection{How function application changes statistical distance}
Let $A$ and $B$ take values in a finite nonempty set $\cX$, and let $\cY$ be a finite nonempty set. Show that for every function $f: \cX \rightarrow \cY$ we have
\[
\SD{f(A)}{f(B)} \leq \SD{A}{B}.
\]
Show that if $f$ is an injection, then
\[
\SD{f(A)}{f(B)} = \SD{A}{B}.
\]
Deduce that if $f$ is a bijection from $\cX$ onto $\cY$, then $d(f(A)) = d(A)$.




\subsection{One-time pad with imperfect randomness}
Let $(\enc,\dec)$ be the one-time pad encryption scheme for messages from $\bin^t$. Suppose a key $K$ is chosen from $\bin^t$ according to some distribution $\alpha$. Consider the following indistinguishability game played between a computationally unbounded adversary $\cA$ and a challenger $\Omega$:
\begin{enumerate}
	\item $\cA$ produces two messages $m_0,m_1 \in \bin^t$ and sends them to $\Omega$.
	\item $\Omega$ selects $b\leftarrow\bin$ uniformly at random, independently samples a fresh key $K$ according to $\alpha$, computes $c = \enc(K,m_b)$, and sends $c$ to $\cA$.
	\item $\cA$ receives $c$ and produces as output $\hat b \in \bin$.
\end{enumerate}
We assume that $\cA$ knows $\alpha$. We say that $\cA$ \emph{won the game} if $\hat b = b$. Show that
\begin{equation*}
\forall_\cA \prob{\cA\mbox{\ wins the game}} \leq \half + d(\alpha).
\end{equation*}

\subsection{Conditional statistical distance}
Let $A$ and $B$ take values in finite nonempty sets $\cX$ and $\cY$, respectively. For each $b$ with $\prob{B=b}>0$, let $\alpha_b$ be the conditional distribution of $A$ given $B=b$, viewed as a distribution over $\cX$. Define the \emph{statistical distance of $A$ from uniformity conditioned on $B$} as
\[
d(A\mid B) := \sum_{b \in \cY:\, \prob{B=b}>0} \prob{B=b} \cdot d(\alpha_b).
\]
Show that $d(A\mid B) = \SD{(A,B)}{(U_\cX,B)}$, where $U_\cX$ is uniform over $\cX$ and independent of $B$.

\noindent
Can you find a game-based interpretation of $d(A\mid B)$ similar to the one in~\ref{ex:game}?

\subsection{Conditioning does not decrease the distance from uniformity}
Show that for every two random variables $A$ and $B$ taking values in finite nonempty sets, we have $d(A\mid B) \geq d(A)$.





\end{document}
